%% LyX 2.4.4 created this file.  For more info, see https://www.lyx.org/.
%% Do not edit unless you really know what you are doing.
\documentclass[english]{article}
\usepackage[T1]{fontenc}
\usepackage[latin9]{inputenc}
\usepackage{float}
\usepackage{enumitem}
\usepackage{amsmath}
\usepackage{amsthm}
\usepackage{amssymb}
\usepackage{geometry}
\geometry{verbose,tmargin=1in,bmargin=1in,lmargin=1in,rmargin=1in}

\makeatletter

%%%%%%%%%%%%%%%%%%%%%%%%%%%%%% LyX specific LaTeX commands.
%% Because html converters don't know tabularnewline
\providecommand{\tabularnewline}{\\}

%%%%%%%%%%%%%%%%%%%%%%%%%%%%%% Textclass specific LaTeX commands.
\numberwithin{equation}{section}
\numberwithin{figure}{section}
\newlength{\lyxlabelwidth}      % auxiliary length 

%%%%%%%%%%%%%%%%%%%%%%%%%%%%%% User specified LaTeX commands.
\usepackage{fancyhdr}
\usepackage{lastpage}
\pagestyle{fancy}
\usepackage{enumitem}
\usepackage{ifthen}
\everymath=\expandafter{\the\everymath\displaystyle}
%\DeclareMathSizes{10}{10}{10}{10}
\usepackage{multicol}

\usepackage{tikz}
\usetikzlibrary{patterns}
\usetikzlibrary{plotmarks}
\usepackage{pgfplots}
\pgfplotsset{compat=newest}

\usepackage{calc}
\usepackage[nomessages]{fp}

\usepackage{datetime}

%%%%%commands
\newcommand{\mb}[1]{\mathbb{#1}}
\newcommand{\funct}[2]{$f\colon#1\rightarrow#2$}
% Added by lyx2lyx
\setlength{\parskip}{\smallskipamount}
\setlength{\parindent}{0pt}

\makeatother

\usepackage{babel}
\begin{document}
\renewcommand{\author}{Dr.\,James Ashe}

\newcommand{\classNum}{439}

\newcommand{\classSec}{A}

\newcommand{\nameType}{solutions}

\newcommand{\examType}{Homework 1}

\newcommand{\Date}{\formatdate{14}{1}{2014}} %%\AdvanceDate[12] \today

%\newdate{my_date}{\day}{\month}{\year}%   
%\AdvanceDate[-5] 
%\displaydate{my_date}

%%First page's header%%%%%%%%%%%%%%%%%%%%%%
\fancypagestyle{firststyle} %%header settings for first pagr
{
	\fancyhead[L]{MTH \classNum{}(\classSec{}) \author{}\\
	\textbf{\examType{}} \Date{}}
	\fancyhead[C]{\textbf{\nameType}}
	\setlength{\headsep}{14pt}
}
%%%%%%%%%%%%%%%%%%%%%%%%%%%%%%%%%%%%%%%%%%

%%Next page's header%%%%%%%%%%%%%%%%%%%%%%%
%\setlist{nolistsep} %eliminates space between list items
\lhead{MTH \classNum{}(\classSec{})} 
\chead{\textbf{\nameType}} 
\cfoot{\thepage{}~of~\pageref{LastPage}} 
\setlength{\headsep}{5pt} 
%%%%%%%%%%%%%%%%%%%%%%%%%%%%%%%%%%%%%%%%%%%

%%Various settings; see the preamble for other settings%%%%%
%%\setenumerate[0]{leftmargin=12pt,itemindent=0pt,labelwidth=0pt}
\renewcommand{\labelenumi}{\textbf{\arabic{enumi}.}}
\renewcommand{\labelenumii}{\textbf{\alph{enumii}.}}
\thispagestyle{firststyle}
%%%%%%%%%%%%%%%%%%%%%%%%%%%%%%%%%%%%%%%%%%

\textbf{Homework}
\begin{enumerate}
\item Show that $\mathbb{Z}_{6}$ is a ring with unity under addition and
multiplication; this covers parts 

\begin{proof}
We know that $\mb{Z}_{6}$ is a commutative group with additive identity
$0$. This covers parts 1-5 of the definition of a ring. Using the
Cayley tables\footnote{A Cayley table shows all the possible products/sums of all a ring's
elements.}:\\
\begin{table}[H]
\begin{centering}
\begin{tabular}{c|c|c|c|c|c|c|}
\multicolumn{1}{c}{$\times$} & \multicolumn{1}{c}{\textbf{0}} & \multicolumn{1}{c}{\textbf{1}} & \multicolumn{1}{c}{\textbf{2}} & \multicolumn{1}{c}{\textbf{3}} & \multicolumn{1}{c}{\textbf{4}} & \multicolumn{1}{c}{\textbf{5}}\tabularnewline
\cline{2-7}
\textbf{0} & 0 & 0 & 0 & 0 & 0 & 0\tabularnewline
\cline{2-7}
\textbf{1} & 0 & 1 & 2 & 3 & 4 & 5\tabularnewline
\cline{2-7}
\textbf{2} & 0 & 2 & 4 & 0 & 2 & 4\tabularnewline
\cline{2-7}
\textbf{3} & 0 & 3 & 0 & 3 & 0 & 3\tabularnewline
\cline{2-7}
\textbf{4} & 0 & 4 & 2 & 0 & 4 & 2\tabularnewline
\cline{2-7}
\textbf{5} & 0 & 5 & 4 & 3 & 2 & 1\tabularnewline
\cline{2-7}
\end{tabular}
\par\end{centering}
\caption{\label{tab:Cayley-tables-for}Cayley tables for $\times$ in $\mb{Z}_{6}$.}
\end{table}

From the table we can see that $\mb{Z}_{6}$ is closed under multiplication.
Association and distribution is preserved in modular arithmetic so
$\mb{Z}_{6}$ is a ring.
\end{proof}
\begin{itemize}
\item Note that we can also see from the table that $\mb{Z}_{6}$ has a
multiplicative identity so it is a ring with unity, but every element
does not have a multiplicative inverse.
\end{itemize}
\item Show that the set $\sqrt{2}\mathbb{Z}=\{x\colon x=m+\sqrt{2}n\mbox{ for some \ensuremath{m,n\in\mathbb{Z}}}\}$
is a ring.

\begin{proof}
Again we know that $\sqrt{2}\mathbb{Z}$ is a commutative group under
addition. So we only need to show that parts 6-8 of the definition
hold. 

(part 6): Let $a,b\in\sqrt{2}\mathbb{Z}$ then $a=m_{1}+n_{1}\sqrt{2}$
and $b=m_{2}+n_{2}\sqrt{2}$ for some $m_{1},m_{2},n_{1},n_{2}\in\mathbb{Z}$.
\begin{eqnarray*}
ab & = & (m_{1}+\sqrt{2}n_{1})(m_{2}+\sqrt{2}n_{2})\\
 & = & m_{1}m_{2}+\sqrt{2}n_{1}m_{2}+\sqrt{2}n_{2}m_{1}+2n_{1}n_{2}\\
 & = & (2n_{1}n_{2}+m_{1}m_{2})+\sqrt{2}(m_{1}n_{2}+m_{2}n_{1})
\end{eqnarray*}
Since integers are closed under addition, $(2n_{1}n_{2}+m_{1}m_{2})$
and $(m_{1}n_{2}+m_{2}n_{1})$ are both integers $ab\in\sqrt{2}\mathbb{Z}$.

Parts 7 and 8 hold since any number in $\sqrt{2}\mathbb{Z}$ is also
a real number, and distribution and association hold in $\mb{R}$.
So $\sqrt{2}\mathbb{Z}$ is a ring. 
\end{proof}
\item Let $F$ be the set of all functions from $\mathbb{R}$ to $\mathbb{R}$.
Define addition and multiplication as 
\[
(f+g)(x)=f(x)+g(x)\mbox{ and }(fg)(x)=f(x)g(x).
\]
Show that $F$ is a commutative ring with identity. 

\begin{proof}
Parts 1 and 6. Let $f,g\in F$. Because $f$ and $g$ are both functions,
each $x$ is mapped to two and only two real numbers: $f(x)$ and
$g(x)$ ($f$ and $g$ are well defined). The operations of addition
and multiplication are also well defined so then $f(x)+g(x)$ and
$f(x)g(x)$ each map $x$ to one real number. So then $(f+g)(x)$
and $(fg)(x)$ are functions in $F$.

Part 8. Let $f,g,h\in F$. $(h(f+g))(x)=h(x)(f+g)(x)=h(x)(f(x)+g(x))$.
So distribution holds.

Parts 2, 3, and 7 hold because they hold in $\mathbb{R}$.

Parts 4 and 5. $0_{F}$ is the zero function $f(x)=0$. $g(x)=-1$
is also a function over $\mathbb{R}$; so then the additive inverse
of $f(x)$ is just $-1f(x)$.

Finally $f(x)=1$ is the multiplicative inverse. So $F$ is a ring
with identity.
\end{proof}
\item Prove that the ring above, $F$, is not an integral domain.

\begin{itemize}
\item [Hint]Here are two \emph{indicator functions} 
\[
\mathbf{1}_{[0,\infty)}(x)=\begin{cases}
1 & \mbox{if }x\in[0,\infty)\\
0 & \mbox{if }x\notin[0,\infty)
\end{cases}\mbox{ and }\mathbf{1}_{(-\infty,0]}(x)=\begin{cases}
1 & \mbox{if }x\in(-\infty,0]\\
0 & \mbox{if }x\notin(-\infty,0]
\end{cases}
\]
\\
 The indicator function is a function that has an output of either
$1$ or $0$. For example, $\mathbf{1}_{[0,\infty)}(2)=1$ and $\mathbf{1}_{[0,\infty)}(-2)=1$.
Now consider the functions $f(x)=x\mathbf{1}_{[0,\infty)}$ and $g(x)=x\mathbf{1}_{(-\infty,0]}$.
\end{itemize}
\begin{proof}
Following the hint, $f(x)\neq0_{F}$ and $g(x)\neq0_{F}$ but $f(x)g(x)=x\mathbf{1}_{[0,\infty)}x\mathbf{1}_{(-\infty,0]}=x^{2}\mathbf{1}_{[0,\infty)}\mathbf{1}_{(-\infty,0]}=0$
if $x=0$ then $f(0)g(0)=0\cdot1^{2}=0$ and if $x\neq0$ then $f(x)g(x)=x^{2}0=0$.
So $F$ is not an integral domain.
\end{proof}
\item Let $R$ and $S$ be rings. Define addition and multiplication on
the Cartesian product $R\times S$ by, 
\[
(r,s)+(r',s')=(r+r',s+s')\mbox{ and }(r,s)(r',s')=(rr',ss')
\]
for $r,r'\in R$ and $s,s'\in S$. Prove that $R\times S$ is a ring.

\begin{proof}
Since $R$ and $S$ are rings $r+r',rr'\in R$ and $s+s',ss'\in S$
so then $(r+r',s+s'),(rr',ss')\in R\times S$. Similarly, parts 2,
3, 5, and 7 hold. 

$(r'',s'')((r,s)+(r',s'))=(r''(r+r'),s''(s+s'))$ for element in $R$
and $S$ so then distribution holds.

Finally, $0_{R\times S}=(0_{R},0_{S})$ so $R\times S$ is a ring.
\end{proof}
\item Show that the set $k\mathbb{Z}=\{x\colon x=kn\mbox{ for some \ensuremath{n\in\mathbb{Z}}}\}$
for any integer $k$ is a subring of $\mathbb{Z}$. 

\begin{proof}
Let $n,m\in\mathbb{Z}$ so that $kn,km\in k\mathbb{Z}$. $(kn)(km)=k(knm)\in k\mathbb{Z}$,
and $kn-km=k(n-m)\in k\mathbb{Z}$ since $knm,n-m\in\mathbb{Z}$.
We know that $\mathbb{Z}$ is a ring. So then $k\mathbb{Z}$ is a
subring (and also a ring) by Theorem 1.3 (the subring test).
\end{proof}
\item Prove that $\mathbb{Q}$ is an integral domain.

\begin{proof}
Let $a,b,c,d\in\mathbb{Z}-\{0\}$ so that $\frac{a}{b},\frac{c}{d}$
are nonzero elements in $\mathbb{Q}$. Then $\frac{a}{b}\frac{c}{d}=\frac{ac}{bd}\neq0$
since $ac\ne0$ since $\mathbb{Z}$ is an integral domain. 
\end{proof}
\item Prove that $\mathbb{Z}_{6}$ is not an integral domain.

\begin{proof}
Consider the nonzero elements $2,3\in\mb{Z}_{6}$, and $2\cdot3=0\bmod6$.
Also see table \ref{tab:Cayley-tables-for}.
\end{proof}
\item Let $a$ and $b$ be two nonzero elements in a ring $R$. Prove that
in general, $ax=b$ does not have a unique solution.

\begin{itemize}
\item [Hint]It should be understood that $x$ must be in $R$. An example
is all that is needed. You can find an example that has two solutions,
or find an example with no solution. 
\end{itemize}
\begin{proof}
$\mb{Z}_{6}$ is a ring, but $2\cdot3=0\bmod6$ and $2\cdot0=0\bmod6$.
So the equation $2x=0$ has more than one solution. See table \ref{tab:Cayley-tables-for}
for more examples.
\end{proof}
\item Prove that $\mathbb{Z}_{n}$ is not an integral domain (under the
usual addition and multiplication) if $n$ is not a prime ($n$ is
assumed to be a positive integer). 

\begin{itemize}
\item [Hint]If $n$ is not prime then $n=mq$ for two integers $1<m,q<n$.
Looking at the example seen in exercise might help. 
\end{itemize}
\begin{proof}
Since $n$ is not prime we can write $n=mq$ for two integers $1<m,q<n$.
$1<m<n$ and $1<q<n$ so both are nonzero elements in $\mathbb{Z}_{n}$,
and $mq=n=0\bmod n$. So $\mathbb{Z}_{n}$ is not an integral domain.
\end{proof}
\end{enumerate}
\textbf{Presentation Problems} 
\begin{enumerate}
\item Prove that the intersection of any two subrings of a ring $R$ is
also a subring of $R$. 

\begin{proof}
Let $S,T$ be two subrings of $R$. $S\cap T$ is a nonempty subset
of $R$ since $0_{R}\in S$ and $0_{R}\in U$. Let $a,b\in S\cap T$.
$ab\in S$ and $a-b\in S$ since $S$ is a ring. Similarly $ab\in T$
and $a-b\in T$. So $ab,a-b\in S\cap T$, and $S\cap T$ is a subring
of $R$ by Theorem 1.3. 
\end{proof}
\item Give an example of a ring with elements $a$ and $b$ with properties
that $ab=0$ but $ba\neq0$.

\begin{proof}
Let $\begin{bmatrix}\begin{array}{rr}
1 & 1\\
0 & 0
\end{array}\end{bmatrix},\begin{bmatrix}\begin{array}{rr}
0 & 1\\
0 & 1
\end{array}\end{bmatrix}\in M_{2\times2}(\mathbb{R})$. The zero in this ring is $\begin{bmatrix}\begin{array}{rr}
0 & 0\\
0 & 0
\end{array}\end{bmatrix}$.

$\begin{bmatrix}\begin{array}{rr}
1 & 1\\
0 & 0
\end{array}\end{bmatrix}\begin{bmatrix}\begin{array}{rr}
0 & 1\\
0 & 1
\end{array}\end{bmatrix}=\begin{bmatrix}\begin{array}{rr}
0 & 2\\
0 & 0
\end{array}\end{bmatrix}$ and $\begin{bmatrix}\begin{array}{rr}
0 & 1\\
0 & 1
\end{array}\end{bmatrix}\begin{bmatrix}\begin{array}{rr}
1 & 1\\
0 & 0
\end{array}\end{bmatrix}=\begin{bmatrix}\begin{array}{rr}
0 & 0\\
0 & 0
\end{array}\end{bmatrix}$
\end{proof}
\item Let $m$ and $n$ be positive integers with least common multiple
$k$. Prove that $m\mathbb{Z}\cap n\mathbb{Z}=k\mathbb{Z}$.

\begin{proof}
Since $k=\mbox{lcm}(n,m)$, $k=na=mb$ for some $a,b\in\mathbb{Z}$.
So then $xk=x(na)=x(ma)$ for any integer $x$ and $xk\in m\mathbb{Z}\cap n\mathbb{Z}$.
Thus $m\mathbb{Z}\cap n\mathbb{Z}\supseteq k\mathbb{Z}$. 

Let $x\in m\mathbb{Z}\cap n\mathbb{Z}$, so then $x=ma=nb$ meaning
$x$ is a common multiple of $m$ and $n$. Since $k$ is the least
common multiple, $k\leq x$ and $k|x$.\footnote{This was shown in MTH 335.}
So then $m\mathbb{Z}\cap n\mathbb{Z}\subseteq k\mathbb{Z}$.
\end{proof}
\item Prove that a ring that is cyclic under addition is commutative. 

\begin{proof}
Recall that cyclic under addition means that there is a element $a$
in a ring $R$ such that for any $r\in R$ $a+a+\cdots+a=r.$ Let
$s,t\in R$. $st=(\underset{m}{\underbrace{a+a+\cdots+a}})(\underset{n}{\underbrace{a+a+\cdots+a}})=mna^{2}=(\underset{n}{\underbrace{a+a+\cdots+a}})(\underset{m}{\underbrace{a+a+\cdots+a}})=ts$. 
\end{proof}
\item Prove that if a ring element has a multiplicative inverse, it is unique. 
\item Prove that a field is an integral domain. 
\item Prove that $\mathbb{Z}_{p}$ is an integral domain if $p$ is prime
(you can assume that it is a commutative ring with identity). 
\item Prove that $\mathbb{Z}_{p}$ is a field.
\end{enumerate}

\end{document}
